\subsection{Separable algebraic extensions}

DEFINITION 1. --- Let E be an algebraic extension of K; then E is said to be separable (over K) if every subextension F of E of finite degree over K is an etale algebra over K (V, p.~28, Def. 1).

Let E be an extension of finite degree of K. Since every subalgebra of an etale algebra is etale (V, p.~30, Prop. 3), it comes to the same to suppose that E is a separable extension of K, or that E is an etale algebra over K.

PROPOSITION 1. --- Let E be an algebraic extension of K. If E is separable, then every subextension E' of E is separable. Conversely, if every subextension of finite degree of E is separable then E is separable.

This follows immediately from Definition 1.

PROPOSITION 2. --- For a field K to be perfect it is necessary and sufficient that every algebraic extension of K should be separable.

Suppose first that K is perfect. Since a field is a reduced ring, it follows from Lemma 5 (V, p.~35) that every extension of finite degree of K is an etale algebra over K; therefore every algebraic extension of K is separable.

Suppose now that K is an imperfect field of characteristic \(p \neq 0\) . Let \(\Omega\) be an algebraic closure of K. Since K is imperfect, there exists \(b \in K\) not belonging to \(K^{p}\) ; put \(a = b^{1/p}\) . Then the extension \(\mathrm{K}(a)\) of K is p-radical of finite degree. By V, p.~26, Prop. 3, there exists a single K-homomorphism of \(\mathrm{K}(a)\) into \(\Omega\) , and since \([\mathrm{K}(a):\mathrm{K}] > 1\) , the algebra \(\mathrm{K}(a)\) is not etale over K (V, p.~32, Prop. 4). In other words, the extension \(\mathrm{K}(a)\) of finite degree of K is not separable.

COROLLARY. --- Every algebraic extension of a field of characteristic 0, or of a finite field, is separable.

This follows from V, p.~7, Prop. 5.

